This thesis presents, implements, and compares four ways of representing melodic context. It asks whether additional model complexity improves next-note prediction enough to justify the resources required to obtain that improvement.

The comparison covers four alternatives, ranging from models that summarize note changes with counts and a limited number of states to a neural model that can use a longer context. All four use the same representation, and arrangements of the same work are kept together. The sample comes from PDMX, a public collection of digital scores: 2,933 pieces remained after preparing the 3,000 selected files.

With the full training set, the neural model assigned the highest probability to the test notes, and the work-level comparison confirmed its advantage. The ordering remained unchanged when the training windows, finite-model capacity, and work split were modified. Removing window overlap reduced the neural model's margin; increasing the finite model's capacity improved its result without reaching a saturation point.

At the scale studied, the neural model achieved a clear predictive advantage, but it was not superior in every respect. Fitting it took less time than fitting two probabilistic alternatives and more than fitting the count-based model; specialized hardware also increased its estimated cost under the documented pricing scenario. The prototype ran end to end on the available equipment. As there was no user evaluation or reference annotation for phrases, sections, or harmonic functions, the results compare prediction and cost, not structural quality or social usefulness.
